\section{Held-out --- con số được phép công bố}

Mọi con số ở hai mục trước đo trên \NumNDev{} câu development --- chính tập mà
nhóm đã thử 4 prompt $\times$ 3 mức depth rồi chọn cấu hình điểm cao nhất. Chọn
ra cái tốt nhất trong 12 lần thử rồi báo cáo điểm của nó trên chính tập đã dùng
để chọn là \textbf{lệch lạc quan theo cấu trúc}. Vì vậy chúng không phải số
công bố.

\subsection{Quy trình, và dấu vết kiểm chứng được}

\begin{itemize}\small
  \item Tập held-out \NumNHeldout{} câu / \NumNArticles{} bài được chốt từ đầu
        bằng lấy mẫu theo bài có seed \NumHeldoutSeed{}; danh sách id được băm
        và ghi vào manifest chuẩn bị.
  \item Bản ghi quyết định winner ký duyệt lúc \NumWinnerApprovedAt{}, băm
        \texttt{\NumHoDecisionSha\ldots}.
  \item Run held-out bắt đầu \NumHoStartedAt{} --- \textbf{sau} khi bản ghi
        quyết định đã đóng băng, và tham chiếu đúng băm đó.
  \item Chạy \textbf{đúng một lần}. Coverage \NumNHeldout{}/\NumNHeldout{}, 0 lỗi.
  \item Bốn trong năm băm artifact (quyết định, danh sách id, predictions,
        judge results) kiểm lại khớp. Băm \texttt{retrievals.jsonl} khớp với bản
        lưu trong \texttt{heldout\_trace/}; bản sao thứ hai dưới \texttt{runs/}
        là một lần tuần tự hóa lại và không khớp --- không ảnh hưởng kết quả,
        nhưng ghi lại để khỏi hiểu nhầm về sau.
\end{itemize}

Thứ tự thời gian ở đây chính là bằng chứng: không thể chọn winner sau khi đã
thấy điểm held-out.

\subsection{Kết quả}

\begin{center}\small
\begin{tabularx}{0.94\textwidth}{@{}L p{2.4cm} p{2.6cm} p{2.2cm}@{}}
\toprule
\NumWinner{} & \textbf{Dev (\NumNDev)} & \textbf{Held-out (\NumNHeldout)} & \textbf{macro} \\
\midrule
Answer Correctness & \NumDevAC & \textbf{\NumHoAC} & \NumHoACMacro \\
Exact Match & \NumDevEM & \NumHoEM & --- \\
Token F1 & \NumDevFone & \NumHoFone & \NumHoFoneMacro \\
Faithfulness & \NumDevFaith & \NumHoFaith & \NumHoFaithMacro \\
Citation F1 & \NumDevCitFone & \NumHoCitFone & \NumHoCitFoneMacro \\
Citation Validity & \NumDevCitVal & \NumHoCitVal & \NumHoCitValMacro \\
Answer Relevancy & \NumDevAnsRel & \NumHoAnsRel & \NumHoAnsRelMacro \\
\midrule
\textbf{Hit@5} (truy xuất) & \NumDevHitFive & \NumHoHitFive & --- \\
\bottomrule
\end{tabularx}
\end{center}

Toàn bộ run tốn \$\NumHoCost{}; latency tổng P50 \NumHoLatency{} ms. Micro và
macro lệch nhau dưới $0{,}008$ ở mọi metric, nên không bài báo nào chi phối con
số công bố.

Điểm cần chú ý: Answer Correctness chỉ tụt $0{,}0193$, nhưng Hit@5 tụt
\NumHoHitFiveDrop{}. \textbf{Tập held-out khó hơn hẳn ở tầng truy xuất}, chứ
không phải ở tầng sinh.

\subsection{Phát hiện quan trọng nhất: prompt không bị overfit}

Tách \NumNHeldout{} câu held-out theo cùng tiêu chí đã dùng ở development:

\begin{center}\small
\begin{tabularx}{0.94\textwidth}{@{}L p{1.2cm} p{2.2cm} p{2cm} p{2cm}@{}}
\toprule
Held-out & \textbf{n} & \textbf{Answer Corr.} & \textbf{Citation F1} & \textbf{Faithfulness} \\
\midrule
\texttt{gold\_in\_top5} & \NumHoNHit & \NumHoHitAC & \NumHoHitCitFone & \NumHoHitFaith \\
\texttt{gold\_not\_in\_top5} & \NumHoNMiss & \NumHoMissAC & \NumHoMissCitFone & \NumHoMissFaith \\
\bottomrule
\end{tabularx}
\end{center}

Trên các câu mà truy xuất làm đúng việc, held-out đạt \NumHoHitAC{} với khoảng tin
cậy 95\% là [0{,}7391; 0{,}7978] --- khoảng này chứa giá trị development
(\NumPtwodFiveHitAC{}), nên hai bên không phân biệt được. Tỉ lệ truy xuất trượt thì tăng gần ba
lần: \NumNMiss{}/\NumNDev{} $\rightarrow$ \NumHoNMiss{}/\NumNHeldout{}.

Citation F1 trên nhóm trượt bằng \NumHoMissCitFone{}, đúng theo định nghĩa:
không có đoạn đúng nào trong context để mà trích dẫn, nên recall bằng 0.
Faithfulness trên nhóm này cũng tụt xuống \NumHoMissFaith{} --- khi không có
bằng chứng, model bám vào đoạn sai.

Chênh lệch giữa hai tập đến từ nhóm \texttt{gold\_not\_in\_top5} lớn hơn
(\NumNMiss{}/\NumNDev{} $\rightarrow$ \NumHoNMiss{}/\NumNHeldout{}), không từ
việc prompt kém tổng quát trên dữ liệu chưa từng thấy.

\subsection{Trần của tầng sinh nằm ở tầng truy xuất}

Cùng một hiện tượng đo được ở cả hai tập:

\begin{center}\small
\begin{tabularx}{0.94\textwidth}{@{}L p{3cm} p{3.4cm}@{}}
\toprule
Answer Correctness của \NumWinner{} & \texttt{gold\_in\_top5} & \texttt{gold\_not\_in\_top5} \\
\midrule
Development (\NumNDev{} câu) & \NumPtwodFiveHitAC & \NumPtwodFiveMissAC \\
Held-out (\NumNHeldout{} câu) & \NumHoHitAC & \NumHoMissAC \\
\bottomrule
\end{tabularx}
\end{center}

Với Hit@5 $=$ \NumHoHitFive{} trên held-out, khoảng 12\% số câu không có đường
nào để trả lời đúng, bất kể prompt viết thế nào. Đầu tư tiếp vào prompt
engineering sẽ cho lợi ích giảm dần; đòn bẩy còn lại nằm ở \textbf{query
rewriting} và ở chính tầng truy xuất.

Kết luận này rút ra được \emph{chỉ vì} hai phase dùng chung một vết truy xuất
đóng băng và mọi điểm số đều truy ngược được về từng câu hỏi. Đây là lợi ích
trực tiếp của quyết định thiết kế ở mục 1.
